% GENERATED FILE. Edit canonical lessons, then run npm run generate:books.
% canonical-source-hash: fnv1a64:d4c9dfd1f06413a6
% canonical-lessons: KA-C324-agnisamaka, KA-C324-hottiko, KA-C324-samayaprajne, KA-C324-nandisu, KA-C324-raksisu

\chapter{Fire in the Market}
\label{ch:ka-fire-in-the-market}
\hlchaptermodality{324}

\begin{chapteropening}
\textbf{By the end of this chapter:} \emph{I can say a firefighter, to catch fire, presence of mind, to extinguish and to rescue.}

Complete the last lesson of chapter 324: I can say a firefighter, to catch fire, presence of mind, to extinguish and to rescue.
\end{chapteropening}

\section[\texorpdfstring{\kn{ಅಗ್ನಿಶಾಮಕ} (agniśāmaka)}{ಅಗ್ನಿಶಾಮಕ (agniśāmaka)}]{\kn{ಅಗ್ನಿಶಾಮಕ} (agniśāmaka) — a firefighter}

\label{lesson:KA-C324-agnisamaka}



\begin{quote}
Before the new one: say the Kannada for an investigation, then the Kannada for to suspect.
\end{quote}

\subsection*{You'll want to know: \kn{ಅಗ್ನಿಶಾಮಕ}}
\textbf{\kn{ಅಗ್ನಿಶಾಮಕ}} — \emph{agniśāmaka} — \textquotedblleft{}a firefighter\textquotedblright{}.

\begin{quote}
\kn{ಒಂದು} \kn{ವೇಳೆ} \kn{ಬೆಂಕಿ} \kn{ಹತ್ತಿದರೆ} \kn{ಅಗ್ನಿಶಾಮಕರು} \kn{ತಕ್ಷಣ} \kn{ಬರುತ್ತಾರೆ}.
\end{quote}

\begin{quote}
\emph{ondu vēḷe beṅki hattidare agniśāmakaru takṣaṇa baruttāre.}
\end{quote}

\begin{quote}
If a fire breaks out, the firefighters come at once.
\end{quote}

Plural: \kn{ಅಗ್ನಿಶಾಮಕರು}, \emph{agniśāmakaru}.

The fire brigade is \kn{ಅಗ್ನಿಶಾಮಕ} \kn{ದಳ}.

\begin{cousinweb}
\textbf{\kn{ಅಗ್ನಿಶಾಮಕ}} (\emph{agniśāmaka}) is built from Sanskrit: a modern compound of \emph{agni}, \textquotedblleft{}fire\textquotedblright{} (CDIAL 55), and a noun built on the root \emph{śam}, \textquotedblleft{}to calm, to put out\textquotedblright{} (Monier-Williams): one who puts out fire.
\end{cousinweb}

\begin{grammarlens}[title={What you've built}]
One more word for telling what happened.
\end{grammarlens}

\subsection*{Your turn}
\begin{itemize}
\raggedright
  \item \emph{Say it:} \emph{agniśāmaka}
  \item \emph{Say it:} \emph{agniśāmaka}, once more
  \item \emph{Say it:} \emph{ondu vēḷe beṅki hattidare agniśāmakaru takṣaṇa baruttāre.} — If a fire breaks out, the firefighters come at once.
  \item \emph{Recall:} say the Kannada for an investigation, then the Kannada for to suspect, then say \emph{agniśāmaka} again
\end{itemize}

\begin{tcolorbox}[breakable,colback=teal!4,colframe=teal!35!black,title={Before you move on}]
What does \kn{ಅಗ್ನಿಶಾಮಕ} mean? (\textquotedblleft{}A firefighter\textquotedblright{}.) Say it once more.
\end{tcolorbox}
\section[\texorpdfstring{\kn{ಹೊತ್ತಿಕೊ} (hottiko)}{ಹೊತ್ತಿಕೊ (hottiko)}]{\kn{ಹೊತ್ತಿಕೊ} (hottiko) — to catch fire}

\label{lesson:KA-C324-hottiko}



\begin{quote}
Before the new one: say the Kannada for to suspect, then the Kannada for a firefighter.
\end{quote}

\subsection*{You'll want to know: \kn{ಹೊತ್ತಿಕೊ}}
\textbf{\kn{ಹೊತ್ತಿಕೊ}} — \emph{hottiko} — \textquotedblleft{}to catch fire\textquotedblright{}.

\begin{quote}
\kn{ಅಡುಗೆಮನೆಯಲ್ಲಿ} \kn{ಎಣ್ಣೆ} \kn{ಹೊತ್ತಿಕೊಂಡಿತು}.
\end{quote}

\begin{quote}
\emph{aḍugemaneyalli eṇṇe hottikoṇḍitu.}
\end{quote}

\begin{quote}
The oil caught fire in the kitchen.
\end{quote}

Past (it …): \kn{ಹೊತ್ತಿಕೊಂಡಿತು}, \emph{hottikoṇḍitu}.

The thing that burns is the subject: \kn{ಮನೆ} \kn{ಹೊತ್ತಿಕೊಂಡಿತು}, the house caught fire.

\begin{cousinweb}
\textbf{\kn{ಹೊತ್ತಿಕೊ}} (\emph{hottiko}) is native Dravidian: DEDR 4517 lists Kannada \emph{pottu}, \textquotedblleft{}to be kindled, catch fire, flame\textquotedblright{}, and modern Kannada adds -\emph{ko}, from \emph{koḷ} \textquotedblleft{}to take\textquotedblright{} (DEDR 2151).
\end{cousinweb}

\begin{grammarlens}[title={What you've built}]
One more word for telling what happened.
\end{grammarlens}

\subsection*{Your turn}
\begin{itemize}
\raggedright
  \item \emph{Say it:} \emph{hottiko}
  \item \emph{Say it:} \emph{hottiko}, once more
  \item \emph{Say it:} \emph{aḍugemaneyalli eṇṇe hottikoṇḍitu.} — The oil caught fire in the kitchen.
  \item \emph{Recall:} say the Kannada for to suspect, then the Kannada for a firefighter, then say \emph{hottiko} again
\end{itemize}

\begin{tcolorbox}[breakable,colback=teal!4,colframe=teal!35!black,title={Before you move on}]
What does \kn{ಹೊತ್ತಿಕೊ} mean? (\textquotedblleft{}To catch fire\textquotedblright{}.) Say it once more.
\end{tcolorbox}
\section[\texorpdfstring{\kn{ಸಮಯಪ್ರಜ್ಞೆ} (samayaprajñe)}{ಸಮಯಪ್ರಜ್ಞೆ (samayaprajñe)}]{\kn{ಸಮಯಪ್ರಜ್ಞೆ} (samayaprajñe) — presence of mind}

\label{lesson:KA-C324-samayaprajne}



\begin{quote}
Before the new one: say the Kannada for a firefighter, then the Kannada for to catch fire.
\end{quote}

\subsection*{You'll want to know: \kn{ಸಮಯಪ್ರಜ್ಞೆ}}
\textbf{\kn{ಸಮಯಪ್ರಜ್ಞೆ}} — \emph{samayaprajñe} — \textquotedblleft{}presence of mind\textquotedblright{}.

\begin{quote}
\kn{ಹುಡುಗಿಯ} \kn{ಸಮಯಪ್ರಜ್ಞೆಯಿಂದ} \kn{ಎಲ್ಲರೂ} \kn{ಬೆಂಕಿಯಿಂದ} \kn{ತಪ್ಪಿಸಿಕೊಂಡರು}.
\end{quote}

\begin{quote}
\emph{huḍugiya samayaprajñeyinda ellarū beṅkiyinda tappisikoṇḍaru.}
\end{quote}

\begin{quote}
Thanks to the girl's presence of mind, everyone escaped the fire.
\end{quote}

Quick, calm thinking in danger: \kn{ಸಮಯಪ್ರಜ್ಞೆ} \kn{ತೋರಿಸು}, to show presence of mind.

\begin{cousinweb}
\textbf{\kn{ಸಮಯಪ್ರಜ್ಞೆ}} (\emph{samayaprajñe}) is borrowed from Sanskrit, a compound of \emph{samaya}, \textquotedblleft{}a meeting, an occasion, time\textquotedblright{} (CDIAL 13185), and \emph{prajñā}, \textquotedblleft{}wisdom, awareness\textquotedblright{} (CDIAL 8515).
\end{cousinweb}

\begin{grammarlens}[title={What you've built}]
One more word for telling what happened.
\end{grammarlens}

\subsection*{Your turn}
\begin{itemize}
\raggedright
  \item \emph{Say it:} \emph{samayaprajñe}
  \item \emph{Say it:} \emph{samayaprajñe}, once more
  \item \emph{Say it:} \emph{huḍugiya samayaprajñeyinda ellarū beṅkiyinda tappisikoṇḍaru.} — Thanks to the girl's presence of mind, everyone escaped the fire.
  \item \emph{Recall:} say the Kannada for a firefighter, then the Kannada for to catch fire, then say \emph{samayaprajñe} again
\end{itemize}

\begin{tcolorbox}[breakable,colback=teal!4,colframe=teal!35!black,title={Before you move on}]
What does \kn{ಸಮಯಪ್ರಜ್ಞೆ} mean? (\textquotedblleft{}Presence of mind\textquotedblright{}.) Say it once more.
\end{tcolorbox}
\section[\texorpdfstring{\kn{ನಂದಿಸು} (nandisu)}{ನಂದಿಸು (nandisu)}]{\kn{ನಂದಿಸು} (nandisu) — to extinguish; to put out}

\label{lesson:KA-C324-nandisu}



\begin{quote}
Before the new one: say the Kannada for to catch fire, then the Kannada for presence of mind.
\end{quote}

\subsection*{You'll want to know: \kn{ನಂದಿಸು}}
\textbf{\kn{ನಂದಿಸು}} — \emph{nandisu} — \textquotedblleft{}to extinguish; to put out\textquotedblright{}.

\begin{quote}
\kn{ಅಗ್ನಿಶಾಮಕರು} \kn{ಒಂದು} \kn{ಗಂಟೆಯಲ್ಲಿ} \kn{ಬೆಂಕಿಯನ್ನು} \kn{ನಂದಿಸಿದರು}.
\end{quote}

\begin{quote}
\emph{agniśāmakaru ondu gaṇṭeyalli beṅkiyannu nandisidaru.}
\end{quote}

\begin{quote}
The firefighters put out the fire in an hour.
\end{quote}

Past (I …): \kn{ನಂದಿಸಿದೆ}, \emph{nandiside}.

Fire and a lamp are put out: \kn{ದೀಪವನ್ನು} \kn{ನಂದಿಸು}. A fire going out by itself is \kn{ನಂದು}.

\begin{cousinweb}
\textbf{\kn{ನಂದಿಸು}} (\emph{nandisu}) is native Dravidian: the causative in -\emph{isu} of \emph{nandu}, \textquotedblleft{}to be quenched, to go out\textquotedblright{}, which DEDR 3590 lists for Kannada.
\end{cousinweb}

\begin{grammarlens}[title={What you've built}]
One more word for telling what happened.
\end{grammarlens}

\subsection*{Your turn}
\begin{itemize}
\raggedright
  \item \emph{Say it:} \emph{nandisu}
  \item \emph{Say it:} \emph{nandisu}, once more
  \item \emph{Say it:} \emph{agniśāmakaru ondu gaṇṭeyalli beṅkiyannu nandisidaru.} — The firefighters put out the fire in an hour.
  \item \emph{Recall:} say the Kannada for to catch fire, then the Kannada for presence of mind, then say \emph{nandisu} again
\end{itemize}

\begin{tcolorbox}[breakable,colback=teal!4,colframe=teal!35!black,title={Before you move on}]
What does \kn{ನಂದಿಸು} mean? (\textquotedblleft{}To extinguish; to put out\textquotedblright{}.) Say it once more.
\end{tcolorbox}
\section[\texorpdfstring{\kn{ರಕ್ಷಿಸು} (rakṣisu)}{ರಕ್ಷಿಸು (rakṣisu)}]{\kn{ರಕ್ಷಿಸು} (rakṣisu) — to rescue; to protect}

\label{lesson:KA-C324-raksisu}



\begin{quote}
Before the new one: say the Kannada for presence of mind, then the Kannada for to extinguish.
\end{quote}

\subsection*{You'll want to know: \kn{ರಕ್ಷಿಸು}}
\textbf{\kn{ರಕ್ಷಿಸು}} — \emph{rakṣisu} — \textquotedblleft{}to rescue; to protect\textquotedblright{}.

\begin{quote}
\kn{ಅಗ್ನಿಶಾಮಕರು} \kn{ಆ} \kn{ಮನೆಯಿಂದ} \kn{ಮೂರು} \kn{ಮಕ್ಕಳನ್ನು} \kn{ರಕ್ಷಿಸಿದರು}.
\end{quote}

\begin{quote}
\emph{agniśāmakaru ā maneyinda mūru makkaḷannu rakṣisidaru.}
\end{quote}

\begin{quote}
The firefighters rescued three children from that house.
\end{quote}

Past (I …): \kn{ರಕ್ಷಿಸಿದೆ}, \emph{rakṣiside}.

Who is rescued takes -\kn{ಅನ್ನು}; what from takes -\kn{ಇಂದ}.

\begin{cousinweb}
\textbf{\kn{ರಕ್ಷಿಸು}} (\emph{rakṣisu}) is built on Sanskrit \emph{rakṣaṇa}, \textquotedblleft{}guarding, tending, protection\textquotedblright{} (CDIAL 10546), with the Kannada verb ending -\emph{isu}.
\end{cousinweb}

\begin{grammarlens}[title={What you've built}]
One more word for telling what happened.
\end{grammarlens}

\subsection*{Your turn}
\begin{itemize}
\raggedright
  \item \emph{Say it:} \emph{rakṣisu}
  \item \emph{Say it:} \emph{rakṣisu}, once more
  \item \emph{Say it:} \emph{agniśāmakaru ā maneyinda mūru makkaḷannu rakṣisidaru.} — The firefighters rescued three children from that house.
  \item \emph{Recall:} say the Kannada for presence of mind, then the Kannada for to extinguish, then say \emph{rakṣisu} again
\end{itemize}

\begin{tcolorbox}[breakable,colback=teal!4,colframe=teal!35!black,title={Before you move on}]
What does \kn{ರಕ್ಷಿಸು} mean? (\textquotedblleft{}To rescue; to protect\textquotedblright{}.) Say it once more.
\end{tcolorbox}
